
\documentclass[10pt,a4paper]{article}

% --- Pakiety ---
\usepackage[utf8]{inputenc}
\usepackage[T1]{fontenc}
\usepackage[polish]{babel}
\usepackage[a4paper,margin=1.5cm]{geometry}
\usepackage{xcolor}
\usepackage{enumitem}
\usepackage{titlesec}
\usepackage{hyperref}
\usepackage{tabularx}
\usepackage{lmodern}

% --- Kolory ---
\definecolor{maincolor}{RGB}{30,65,100}
\definecolor{graytext}{RGB}{90,90,90}

% --- Formatowanie ---
\pagestyle{empty}
\setlength{\parindent}{0pt}
\setlength{\parskip}{3pt}
\setlength{\tabcolsep}{0pt}

\titleformat{\section}
{\large\bfseries\color{maincolor}}
{}{0em}{}[\titlerule]

\titlespacing{\section}{0pt}{9pt}{5pt}

\setlist[itemize]{
    leftmargin=15pt,
    itemsep=2pt,
    topsep=3pt
}

\hypersetup{
    colorlinks=true,
    urlcolor=maincolor
}

\begin{document}

% --- Dane osobowe ---
\begin{center}
    {\Huge \textbf{\color{maincolor} KRYSTIAN CZARNOTA}}\\[5pt]
    {\large Programista}\\[7pt]

    Wrocław, Polska \quad | \quad
    +48 123 456 789
\end{center}

% --- Profil ---
\section*{Profil zawodowy}

Programista zainteresowany tworzeniem oprogramowania,
aplikacji oraz rozwiązań internetowych. Posiadam doświadczenie
w pracy z językami Java, JavaScript i C\#, a także w tworzeniu
projektów z wykorzystaniem silnika Unity. Interesuję się
również web designem, projektowaniem interfejsów oraz
tworzeniem nowoczesnych i funkcjonalnych stron internetowych.
Jestem osobą szybko uczącą się, nastawioną na rozwój
i chętną do zdobywania nowych umiejętności.

% --- Doświadczenie i projekty ---
\section*{Doświadczenie i projekty programistyczne}

\textbf{Tworzenie aplikacji i oprogramowania}
\begin{itemize}
    \item Rozwijanie aplikacji z wykorzystaniem języków Java
    oraz JavaScript.
    \item Implementowanie funkcjonalności i rozwiązywanie
    problemów programistycznych.
    \item Praca z kodem oraz rozwijanie umiejętności
    projektowania oprogramowania.
\end{itemize}

\textbf{Tworzenie gier i aplikacji w Unity}
\begin{itemize}
    \item Tworzenie projektów z wykorzystaniem języka C\#
    i silnika Unity.
    \item Implementowanie mechanik oraz funkcjonalności
    interaktywnych.
\end{itemize}

\textbf{Web Design i tworzenie stron internetowych}
\begin{itemize}
    \item Projektowanie nowoczesnych i estetycznych
    interfejsów stron internetowych.
    \item Dbanie o funkcjonalność, czytelność
    i wygodę użytkowania.
\end{itemize}

% --- Wykształcenie ---
\section*{Wykształcenie}

\textbf{Uniwersytet Wrocławski}
\hfill \textcolor{graytext}{2026 -- obecnie}\\
Informatyka -- student

% --- Umiejętności ---
\section*{Umiejętności}

\begin{tabularx}{\textwidth}{@{}l X@{}}
\textbf{Języki programowania:} &
Java, JavaScript, C\#, Python, SQL \\[3pt]

\textbf{Technologie:} &
Unity, Git, Linux \\[3pt]

\textbf{Web Design:} &
Projektowanie stron internetowych,
projektowanie interfejsów UI/UX \\[3pt]

\textbf{Narzędzia:} &
Visual Studio Code, Excel \\[3pt]

\textbf{Języki obce:} &
Polski -- ojczysty, Angielski -- C2 \\[3pt]

\textbf{Umiejętności miękkie:} &
Analityczne myślenie, komunikatywność,
praca zespołowa, szybkie przyswajanie wiedzy
\end{tabularx}

% --- Dodatkowe informacje ---
\section*{Dodatkowe informacje}

Zainteresowanie nowymi technologiami, rozwojem
oprogramowania, tworzeniem gier komputerowych
oraz projektowaniem stron internetowych.

\vfill

% --- Klauzula ---
{\footnotesize\color{graytext}
Wyrażam zgodę na przetwarzanie moich danych osobowych
zawartych w CV na potrzeby obecnego procesu rekrutacyjnego,
zgodnie z obowiązującymi przepisami prawa.
}

\end{document}
